%! AP Statistics — Unit 4, Lesson 4.8 — Guided Notes (Answer Key)
\documentclass[10pt]{article}
\usepackage{apstats-article}
\usepackage{apstats-key}
\newcommand{\TallMath}[1]{$\displaystyle #1\rule[-1.4em]{0pt}{3.2em}$}

\begin{document}

\pageheader{Unit 4, Lesson 4.8}{Guided Notes --- Answer Key}
\namedateperiod

\vspace{0.12in}

\begin{objectivebox}
\small By the end of this lesson, I will be able to\dots\\[4pt]
\ansline{Compute the expected value: $\mu_X = \sum x \cdot P(X=x)$.}\\[1pt]
\ansline{Compute the standard deviation $\sigma_X$ of a discrete random variable.}\\[1pt]
\ansline{Interpret $\mu_X$ and $\sigma_X$ in context.}
\end{objectivebox}

\vspace{0.1in}

\begin{vocabbox}
\noindent\textbf{\textcolor{navy}{Expected value ($\mu_X$):}}\\[1pt]%
\ansline{The long-run average value of $X$: a weighted average of its possible values.}\\[1pt]\ansline{}\\[3pt]
\noindent\textbf{\textcolor{navy}{Variance ($\sigma^2_X$):}}\\[1pt]%
\ansline{The probability-weighted average squared deviation from the mean:}\\[1pt]\ansline{$\sigma^2_X = \sum (x-\mu_X)^2 \cdot P(X=x)$.}\\[3pt]
\noindent\textbf{\textcolor{navy}{Standard deviation ($\sigma_X$):}}\\[1pt]%
\ansline{$\sigma_X = \sqrt{\sigma^2_X}$: the typical distance of $X$ from $\mu_X$, in the units of $X$.}\\[1pt]\ansline{}\\[3pt]
\noindent\textbf{\textcolor{navy}{Long-run average:}}\\[1pt]%
\ansline{The value the mean of the outcomes approaches over many repetitions.}\\[1pt]\ansline{}\\[3pt]
\end{vocabbox}

\vspace{0.1in}

\begin{hookbox}
\small
A carnival game costs \$2 to play. Prize distribution:
\begin{center}
\renewcommand{\arraystretch}{1.3}
\begin{tabular}{@{}ccc@{}}
Prize ($x$) & \$0 & \$5 \\
\hline
Probability & 0.70 & 0.30 \\
\end{tabular}
\end{center}
If you play 100 times, what do you expect to win in total? \ans{$100(\$1.50) = \$150$}\\
Is playing worth the \$2 entry fee? \ans{No: \$1.50 per play $<$ \$2}
\end{hookbox}

\vspace{0.1in}

\begin{notesbox}{Mean (Expected Value) of a Random Variable}
\small

The \textbf{expected value} of $X$ is a \ans{weighted} average of all possible values,
weighted by their probabilities:

\vspace{8pt}
\begin{center}
\TallMath{\mu_X = \sum x \cdot P(X=x) = x_1 \cdot P(X=x_1) + x_2 \cdot P(X=x_2) + \cdots}
\end{center}

\vspace{8pt}
\textbf{Computation Table Example:} $X$ = prize from carnival game.

\vspace{4pt}
\renewcommand{\arraystretch}{1.8}
\begin{tabularx}{\linewidth}{@{} c c X @{}}
$x$ & $P(X=x)$ & $x \cdot P(X=x)$ \\
\hline
\$0 & 0.70 & \ans{$0(0.70) = \$0.00$} \\
\$5 & 0.30 & \ans{$5(0.30) = \$1.50$} \\
\end{tabularx}

\vspace{6pt}
$\mu_X = \sum x \cdot P(X=x) = $ \ans{$\$0 + \$1.50 = \$1.50$}

\vspace{6pt}
\textbf{Interpret} $\mu_X$ in context: \ans{$\mu_X = \$1.50$} means that in the \ans{long run (many plays)},\\
the average prize per play is approximately \ans{\$1.50}.

\end{notesbox}

\vspace{0.1in}

\begin{notesbox}{Standard Deviation of a Random Variable}
\small

The \textbf{variance} measures how spread out the values of $X$ are:

\vspace{6pt}
\begin{center}
\TallMath{\sigma^2_X = \sum (x - \mu_X)^2 \cdot P(X=x)}
\end{center}

\vspace{6pt}
The \textbf{standard deviation} is $\sigma_X = $ \ans{$\sqrt{\sigma^2_X}$}.

\vspace{8pt}
\textbf{Full Computation Table:} $X$ = number of cars per household ($\mu_X = 1.55$).

\vspace{4pt}
\renewcommand{\arraystretch}{2.0}
\begin{tabularx}{\linewidth}{@{} c c c c X @{}}
$x$ & $P(X=x)$ & $x \cdot P(X=x)$ & $(x - \mu_X)^2$ & $(x-\mu_X)^2 \cdot P(X=x)$ \\
\hline
0 & 0.10 & \ans{0.00} & \ans{2.4025} & \ans{0.2403} \\
1 & 0.40 & \ans{0.40} & \ans{0.3025} & \ans{0.1210} \\
2 & 0.35 & \ans{0.70} & \ans{0.2025} & \ans{0.0709} \\
3 & 0.15 & \ans{0.45} & \ans{2.1025} & \ans{0.3154} \\
\hline
\textbf{Sum} & & $\mu_X =$ \ans{1.55} & & $\sigma^2_X =$ \ans{0.7475} \\
\end{tabularx}

\vspace{6pt}
$\sigma_X = \sqrt{\sigma^2_X} = $ \ans{$\sqrt{0.7475} \approx 0.865$ cars}

\vspace{6pt}
\textbf{Interpret} $\sigma_X$ in context:\\
\ansline{Cars per household typically vary from the mean of 1.55 cars by about 0.86 cars.}

\end{notesbox}

\vspace{0.1in}

\begin{practicebox}
\small
$X$ = number of goals scored per soccer game.
\begin{center}
\renewcommand{\arraystretch}{1.4}
\begin{tabular}{@{}cccc@{}}
$x$ & 0 & 1 & 2 \\
\hline
$P(X=x)$ & 0.30 & 0.50 & 0.20 \\
\end{tabular}
\end{center}

\renewcommand{\arraystretch}{2.0}
\begin{tabularx}{\linewidth}{@{} c c c c X @{}}
$x$ & $P(X=x)$ & $x \cdot P(X=x)$ & $(x - \mu_X)^2$ & $(x-\mu_X)^2 \cdot P(X=x)$ \\
\hline
0 & 0.30 & \ans{0.00} & \ans{0.81} & \ans{0.243} \\
1 & 0.50 & \ans{0.50} & \ans{0.01} & \ans{0.005} \\
2 & 0.20 & \ans{0.40} & \ans{1.21} & \ans{0.242} \\
\hline
\textbf{Sum} & & $\mu_X =$ \ans{0.90} & & $\sigma^2_X =$ \ans{0.490} \\
\end{tabularx}

\vspace{6pt}
$\sigma_X = $ \ans{$\sqrt{0.490} \approx 0.70$ goals}

\vspace{6pt}
Interpret $\mu_X$ in context: \ansline{Over many games, the team averages 0.90 goals per game.}

\vspace{2pt}
Interpret $\sigma_X$ in context: \ansline{Goals per game typically vary from 0.90 by about 0.70 goals.}
\end{practicebox}

\end{document}
